\section{Why this proof fits iPRA}\label{sec:arithmetic}

There are two distinct obligations. The output procedures must be
primitive recursive, and their correctness must be provable with the
induction permitted in iPRA. The first alone does not imply the
second. We give the bookkeeping for both.

\subsection{Coding and the exact logical form}

Use a signed-integer code, positive-denominator rational pairs,
and a primitive recursive pairing function. A list can be a length
together with iterated pairs; lookup is bounded iteration of the
projection. For lists whose entries are below $M$, one may instead
use base $M+1$ digits and a length. Lists of length at most $M$
then have a primitive recursive code bound, for example a pairing
of $M$ with $(M+1)^{M+1}$. Thus even a phrase quantifying over
all paths of length at most $M$ is a bounded arithmetic phrase
after this coding.

Rational arithmetic uses cross multiplication and signed arithmetic.
Normalization by a gcd is optional; equality can always be tested by
cross multiplication. If a total routine encounters an invalid code
or a zero denominator outside its intended domain, give it a fixed
default value. Correctness is conditional on the explicit domain
predicates. Division with remainder, bounded minimum, maximum,
finite xor, and finite sorting all have standard bounded-loop
definitions using primitive recursion.

Write $\Reg_G(\vec a)$ for the universal arithmetic presentation
conditions of Section~\ref{sec:names}, with the primitive recursive
definitions fixed. For fixed point names $X,Y$, write
$\operatorname{Off}(X,k)$ for the universal clearance formula
there. The classification is a numerical primitive recursive term
$\operatorname{Bit}(\vec a,\vec x,k)$.
The router is another such term that returns a tag, a finite route
code, the exponent $H$ of its clearance, and finite parent and
distance certificates. Its output terms have total default branches.
The useful assertion has the form
\begin{gather*}
 \Reg_G(\vec a)\wedge
 \operatorname{Off}(X,k_x)\wedge\operatorname{Off}(Y,k_y)
 \\[1mm]
 \longrightarrow\ \bigl(\operatorname{Bit}(X,k_x)=\operatorname{Bit}(Y,k_y)
       \\[-1mm]\hspace{18mm}\longrightarrow\
       \operatorname{RouteCorrect}
          (X,Y,\operatorname{Route}(X,Y,k_x,k_y))\bigr).
\end{gather*}
This is an ordinary first-order formula after expanding the
notations. A route's rational intermediate vertices are finite codes.
Its first and last vertices use the fixed input names.
Their segment interpolations are primitive recursive combinations
of these names and rational coefficients, rather than calls to an
interpreter for arbitrary programs.

The finite part of route correctness can be particularly concrete.
The stored curve net has error at most $h/8$, and every used
center link $S$ satisfies the rational check
\[
                \min_j\dist(S,Y_j)\ge3h/2.
\]
This finite check, the net correctness theorem, and the short-endpoint
inequalities imply the universal clearance conclusion. The output
does not consist merely of a search claim; its finite route and
the data certifying each link are the numerical witnesses.

For the local-boundary procedure the output is a pair of rational
point codes, two finite parity calculations, and the common
clearance exponent $K+4$. Its correctness is an analogous formula
from $\Reg_G$ to the two bit values, the dyadic distance bounds,
and the universal clearance assertions.

Uniformity is by substitution of fixed primitive recursive
definitions into these constructions. A family with additional
number parameters is still one schema instance. We make no
claim that validity of an arbitrary presentation is decidable,
or that iPRA proves the regularity premises of every true
presentation.

\subsection{Padding from rational conditions to completed names}
\label{sec:padding}

To remove a minor convention issue about endpoints of modulus
inequalities, one can begin with rational-parameter moduli
$U_0,V_0$ satisfying \eqref{eq:arithU}--\eqref{eq:arithV}
on rational parameters. Normalize them by finite maxima.
For the completed names use the more generous moduli
\begin{equation}
        U(k)=U_0(k+2)+2,\qquad V(k)=V_0(k+2)+2,
 \label{eq:padded}
\end{equation}
and normalize again if needed.

For example, if completed parameters have distance
at most $2^{-U(k)}$, sufficiently accurate rational approximations
have distance strictly below $2^{-U_0(k+2)}$. Their images
are within $2^{-(k+2)}$ up to the specified approximation
errors, hence within $2^{-k}$ in the limit-name inequalities.
If the completed parameters have distance at least $2^{-k}$,
their sufficiently accurate rational approximations have distance
at least $2^{-(k+2)}$. Their image separation is at least
$2^{-V_0(k+2)}$ up to arbitrarily small errors, and therefore
at least $2^{-V(k)}$.

This passage is an estimate at each requested precision.
It is not an application of compactness or of real-number
trichotomy. All the proof's named constants refer to the padded
$U,V$ if the input was originally given only in rational form.

Evaluation on a completed parameter is explicit as well.
If $q_j$ is a fixed fast circle name for $t$ with
$d_\T(q_j,t)\le2^{-j}$, use
\[
                    G(q_{U(k+2)},k+2)
\]
as an approximation to $\gamma(t)$ at target precision $k$.
Forward continuity and evaluation contribute at most
$2^{-(k+2)}$ each. This is a primitive recursive composition
for the fixed name; an unnamed extension operator is unnecessary.

\subsection{An explicit bound ledger}
\label{sec:bounds}

The following are deliberately generous bounds. They are intended
to demonstrate primitive recursiveness, not practical efficiency.
Let the routing parameters be those of
\eqref{eq:route-scale}--\eqref{eq:board}, and put
$N=2^{U(H+4)+1}$.

\begin{center}
\begin{tabular}{@{}p{0.59\textwidth}p{0.33\textwidth}@{}}
\toprule
Operation & A priori bound\\
\midrule
Curve samples used for link tests & $N$\\
Rational point-to-segment minimization & at most $8$ candidates\\
Board cells & $M=(2L+1)^2$\\
Neighbor-link tests, allowing duplicates & $6M$\\
Reachability propagation & $M$ rounds\\
Parent route in the board & $M-1$ links\\
Route with original endpoint connectors & $M+1$ segments\\
Separating-cycle trace & $6M$ edges\\
Expanded cyclic-mesh walk in factorization & $6MN$ edges\\
Edge-usage bit table & $N$ entries\\
\bottomrule
\end{tabular}
\end{center}

A direct implementation can use $M$ rounds, each scanning all
$M$ cells and six neighbors. This gives at most $6M^2$
Boolean propagation tests after the geometric link table has
been calculated. A queue implementation improves this bound,
but the proof does not need the improvement.
The shortest-coordinate path used to find a separating edge
has at most $4L$ links.

For the local-square argument, use $m,K,T$ from
\eqref{eq:localsquare}--\eqref{eq:Tsquare}. There are
$4T$ segments and no more than $4T$ transition vertices.
One half-arc net with at most
$2^{U(K+9)+1}+1$ samples is sufficient for all color tests.
Every transition vertex has clearance $2^{-(K+4)}$,
so the hypothetical same-bit routes in the contradiction
can all use
\[
 k_{\mathrm{local}}=K+7,\qquad
 H_{\mathrm{local}}=V(U(K+11))+10
\]
and a common board based, for example, on $B=R+8$.
All square vertices have norm below $R+1$ with the generous
bound \eqref{eq:R}. Thus the finite family of at most $4T$
auxiliary route calculations also has a displayed primitive
recursive bound. The witness-finding procedure itself only
needs to compute the transition bits and select the two values.

For approximate intersection at precision $n$, one may use
$2^{W(n+4)+1}+1$ path samples, where $W$ is the path's forward
modulus, and $2^{U(n+4)+1}$ curve samples. The product of
these two finite numbers bounds the pair search.

Every function in this ledger is obtained from supplied fixed
primitive recursive functions by composition, arithmetic,
exponentiation, finite maxima, and bounded iteration. A call
to $U$ or $V$ at a computed argument is just composition.
No step increases its precision until some undecidable
condition happens to become true.

\subsection{Quantifier-free induction, step by step}
\label{sec:qfind}

For a primitive recursive Boolean predicate $A(i,\vec z)$, define its
Boolean fold
\[
 C_A(0,\vec z)=1,\qquad
 C_A(s+1,\vec z)=
     C_A(s,\vec z)\,[A(s,\vec z)].
\]
The finite universal statement
$\forall i<s\,A(i,\vec z)$ is represented by the equation
$C_A(s,\vec z)=1$. Finite existential statements use the
corresponding or-fold; least witnesses are bounded scans.
Hence the induction invariants below are quantifier-free
equations or Boolean combinations of equations in the
language containing the relevant primitive recursive functions.

\begin{longtable}{@{}p{0.24\textwidth}p{0.69\textwidth}@{}}
\toprule
Construction & Induction invariant and step justification\\
\midrule
\endhead
Finite minima &
The stored candidate is among the processed candidates and has
value no larger than each of them. The new candidate is compared
by decidable rational order. A bounded fold records all preceding
comparisons.\\[3pt]
Ray parity &
The accumulator is the xor of the processed edge contributions.
For a triangle the calculation is a fixed finite case analysis;
for a list of triangles, adding the next triangle cancels its
twice-occurring internal edges.\\[3pt]
Mesh refinement &
The current edge chain is the original chain plus the boundaries
of the triangles already inserted. Each new triangle has its
specified margin from the query, so its parity is zero.\\[3pt]
Reachability &
Every discovered cell has a recorded safe path of length at
most the current round. After round $s$, every safe path of at
most $s$ links ends in the reached set. Path codes of these
lengths have an explicit finite bound.\\[3pt]
Parent extraction &
The current parent has strictly smaller discovery rank.
Each chosen parent link is safe and the current reversed list
has no repeated cell. At most $M$ parent lookups suffice.\\[3pt]
Stabilization &
If a reached set has not stabilized, the next round adds a
new cell. A bounded count of distinct reached cells increases
and cannot exceed $M$. Equivalently, bounded loop deletion
shortens every walk to at most $M-1$ links.\\[3pt]
Boundary traversal &
Orient separating edges with green on the left. This table has
one successor and one predecessor at each used vertex. Its current orbit has no
repetition before its return. The local degree statement is
the eight-case calculation for three cell colors.\\[3pt]
Cycle factorization &
At each sample vertex, the xor of the two incident usage bits
is zero. Induction on the cyclic vertex index makes every
usage bit equal to the initial one. The total parity is the
corresponding finite xor identity.\\[3pt]
Square contradiction &
After $j$ monochromatic runs, the xor of the $j$ loop chains
equals the xor of their routes and their square runs.
The equalities at $c,d$ telescope by adding one run.\\[3pt]
Circle-name lift &
After $j$ increments the lift projects to the selected
circle residue, and the absolute accumulated increment is
bounded by the finite geometric sum. Its bound is proved
by rational arithmetic induction.\\
\bottomrule
\end{longtable}

Some invariants are naturally first expressed as bounded
quantification over arrays or lists. The fold translation just
described is part of the proof, so those do not become new
unbounded induction instances. The natural induction counters
are list lengths, board rounds, or positions in finite sums.

The universal presentation hypotheses remain assumptions in
the deduction context. During a step they are instantiated at
the rational parameters and precisions occurring in that step.
Safe-edge and replacement-triangle margins are recorded as finite
rational certificate predicates. Induction verifies those
predicates; the general net lemma then gives the universal
clearance conclusion with an arbitrary curve parameter held fixed.
One does not induct on the unbounded assertion that all
previous paths avoid the named curve.
We do not induct on the assertion that a real-valued function
is continuous, that an arbitrary point belongs to a component,
or that a formula holds for every real point. For a geometric
estimate with an arbitrary query or precision, hold those
number parameters fixed, perform the finite calculation and
its quantifier-free induction, and generalize afterward.

This also explains why the exact common-boundary converse does
not require stronger induction. For fixed $n,j$ the closeness
of the two selected parameters follows directly from their
two distance bounds and the separation condition. Only the
finite sum of lifted increments uses induction, and its
invariant is a rational inequality.

\subsection{Where intuitionistic reasoning matters}

The proof uses case distinctions only on coded rational
comparisons, finite lists, or bits. In the routing argument,
reachability is one such finite Boolean value. In the local
square argument, the assertions that all segment colors agree
and that all transition bits agree are finite Boolean tests.
In the approximate intersection and boundary converse,
the test $D\ge\eps/2$ is a rational comparison.
Showing one branch contradictory therefore permits selection
of the other branch intuitionistically.

The use of a positive inverse modulus is also precise.
Given two rational parameters, compare their distance with
$2^{-r}$. The greater-or-equal case contradicts the
strictly smaller image-distance estimate, so the strict
parameter-distance bound follows. This is not the
contrapositive of mere negative injectivity.

We have used neither unrestricted excluded middle nor
Markov's principle, neither a fan theorem nor weak
K\"onig's lemma, and neither a general extremum theorem
nor an intermediate value selection principle.
The completed-name language packages explicit rational
estimates; it is not an additional axiom asserting the
existence of arbitrary infinite objects.

\begin{corollary}[First-order iPRA reading]\label{cor:ipra}
For every fixed tuple of primitive recursive definitions
for the presentation and any required point and path names,
the arithmetic versions of Theorem~\ref{thm:main} are
derivable in intuitionistic primitive recursive arithmetic.
Their witness procedures are the primitive recursive
constructions above. The derivations are conditional on
the stated arithmetic regularity and clearance premises.
\end{corollary}

\begin{proof}
Introduce the finitely many required primitive recursive
function symbols with their defining equations. Prove the
bounded arithmetic and finite geometry lemmas by the
quantifier-free invariants just listed. Substitute them
into the proofs of Sections~\ref{sec:parity}--\ref{sec:completion}.
All other logical steps are first-order intuitionistic
introduction and elimination rules. Existential output
claims are introduced with the already specified witness
terms. The arbitrary requested precision is then universally
generalized. The construction uses exactly the defining
equations and quantifier-free induction available in iPRA.
\end{proof}

This corollary is a human proof of formalizability with an
explicit induction audit. It is not a claim that the many
object-theory proof terms have already been generated and
checked by Lean. The next section identifies that remaining
mechanical task and the exact existing repository interface.
